% Pillar-3 (group size) standalone introduction + related work
\section{Introduction}
\label{sec:p3-intro}
GRPO~\citep{shao2024deepseekmath, guo2025deepseekr1} estimates advantages by
sampling a group of $G$ completions per prompt and centering their rewards. The
group size $G$ thus plays two roles at once: it is a variance-reduction knob, like
a Monte-Carlo baseline, and it is a preference-density knob, controlling how many
within-prompt comparisons the update sees. Practitioners often treat larger $G$ as
strictly better, reasoning by analogy to variance reduction, but this conflates the
two roles and ignores compute cost. This paper measures the effect of $G$ and
clarifies which mechanism is doing the work.

The question connects GRPO to preference optimization. Direct Preference
Optimization~\citep{rafailov2024direct} learns from explicit pairwise preferences;
GRPO, by centering rewards within a group, induces \emph{implicit} pairwise
contrasts from a scalar reward. Recent analyses formalize this
correspondence~\citep{wu2025grpo_dpo, liu2026gdpo}, and it reframes $G$ as a control
on how many endogenous preference pairs each prompt contributes---nonzero only for
mixed-reward groups.

We study $G$ as one pillar of \ours{}, a controlled benchmark of 70+ runs across
seven RL libraries and five model families ($0.6$B--$\sim$671B) on GSM8K
\citep{cobbe2021training}, HumanEval~\citep{chen2021codex}, and tool-use tasks. Our
contributions are: (i) an illustrative reconstruction showing the effect of $G$ is non-monotone,
with intermediate sizes often best and no universal optimum, illustrated by a
$G{=}4$ versus $G{=}32$ comparison with near-complete reward retention and an
estimated signal-to-noise ratio that grows at only $\approx 52\%$ of the $\sqrt{G}$ ideal;
(ii) an explicit validation, via observable signatures in logged diagnostics, that the group-centered
update is an all-pairs preference contrast---the ``GRPO is secretly DPO'' bridge;
and (iii) equivalence and difficulty-stratified analyses (TOST) that separate a
variance-reduction reading of the $G$ effect from a preference-density one. We make
no claim that any single $G$ or algorithm is universally superior.

\subsection{Related Work}
\label{sec:p3-related}
Group-relative and critic-free RL objectives have proliferated around
GRPO~\citep{shao2024deepseekmath, ahmadian2024backtobasics, lin2025cppo, gift2025}.
The preference-learning view descends from
DPO~\citep{rafailov2024direct} and its relatives~\citep{azar2024ipo, ethayarajh2024kto, hong2024orpo}, and the GRPO$\leftrightarrow$DPO correspondence has been examined
directly~\citep{wu2025grpo_dpo, liu2026gdpo}. Our contribution is empirical: rather
than derive the bridge in isolation, we validate the all-pairs contrast identity
against logged GRPO diagnostic signatures and connect the group-size effect to preference
density on a fixed benchmark.
